\documentclass[11pt]{article}

\usepackage{series}
\usepackage{amsmath,amssymb}
\usepackage{booktabs}
\usepackage{enumitem}

\newcommand{\cP}{\mathcal P}
\newcommand{\cH}{\mathcal H}

\title{A Projection-First Reframing of Quantum Gravity}
\author{Peter Nero}
\date{Second edition\\July 2026}

\begin{document}
\maketitle

\begin{abstract}
This essay develops a projection-first research perspective on quantum
gravity.  It does not deny that metric perturbations have propagating
helicity-two degrees of freedom, that they can be quantized perturbatively, or
that General Relativity is a predictive low-energy quantum effective field
theory.  Instead, it asks whether the effective quantum geometry seen by an
observer could be the retained response of a deeper quantum source and
selection process.  On that view, quantized matter and source action can place
stress on an effective geometric medium, while projection determines which
geometric distinctions remain observable.  Horizons and singularities then
become candidates for observer-relative or theory-relative limits of an
effective description, not automatic proofs of projection failure.  We
separate this interpretation from established results, compare it with
effective-field-theory, holographic, and background-independent approaches,
and list the mathematical gates a constructive MTT realization must pass:
a Lorentzian principal symbol, a positive physical spin-two sector, BRST or
equivalent gauge consistency, universal stress-energy coupling, controlled
local reduction, and global gluing.  Current MTT work supplies an internal
helicity-two support certificate and conditional reduction machinery, but not
a completed ultraviolet theory, Newton normalization, or full matter
response.  Projection-first quantum gravity is therefore presented as a
testable research program rather than a solved theory.
\end{abstract}

\section*{Revision note for this edition}
\addcontentsline{toc}{section}{Revision note for this edition}
\begin{description}
\item[Supersedes.] \emph{A Projection-First Reframing of Quantum Gravity},
version~1.
\item[Reason.] The first edition treated gravity as categorically
unquantizable, described the metric as mere bookkeeping, portrayed existing
quantum-gravity programs as repeated failures, and promoted horizons and
singularities directly to projection breakdown.  Those claims exceeded both
standard results and the current MTT construction.
\item[Resolution.] This edition distinguishes effective metric degrees of
freedom from a proposed deeper source ontology, recognizes perturbative
quantum GR as a valid low-energy EFT, treats horizon and singularity language
as hypotheses with explicit limits, and turns the proposal into a list of
constructive and falsifiable gates.
\item[Retained content.] The projection-first question remains useful:
noninvertible coarse description can explain why globally distinct states
look identical to restricted observers, and geometry may be an effective
response to selected quantum source data rather than the most fundamental
object.
\item[Remaining boundary.] No current MTT theorem derives the complete
Lorentzian gravitational action, its physical normalization, universal
stress response, or ultraviolet completion from projection alone.
\end{description}

\tableofcontents

\section{Overview and How to Read the Reframing}

\subsection{The central picture}

In plain language, the projection-first idea introduces three levels:
\[
\begin{gathered}
\boxed{\text{upper source dynamics}}
\longrightarrow
\boxed{\text{selection and projection}}
\\[1mm]
\Downarrow
\\[-1mm]
\boxed{\text{effective quantum fields and geometry}} .
\end{gathered}
\]
The first level contains whatever dynamics the underlying theory ultimately
uses.  The second specifies which distinctions survive into a particular
effective description.  The third is the world in which laboratory fields,
clocks, rods, causal order, and gravitational waves are represented.

This separation suggests a possibility, not a conclusion.  Quantized source
action might generate stress, while the effective metric records the coherent
response after projection.  In that interpretation one need not regard the
metric as the most fundamental variable.  Nevertheless, the effective metric
can still possess genuine propagating modes and quantum correlations.  Being
emergent does not mean being nondynamical or classical.

\subsection{What the paper does not claim}

The essay does not claim that:
\begin{itemize}
\item quantizing metric perturbations is a category error;
\item horizons locally destroy physical description;
\item every singularity has divergent curvature;
\item area laws prove a finite MTT projection capacity;
\item existing quantum-gravity programs have failed mathematically; or
\item MTT already supplies a complete ultraviolet theory of gravity.
\end{itemize}
Those stronger statements are neither needed nor currently justified.  The
purpose is to identify a different division of labor between quantum source
dynamics, projection, and effective geometry, then ask what would make that
division mathematically credible.

\section{What Established Quantum Gravity Already Tells Us}

\subsection{The metric has physical perturbative degrees of freedom}

General Relativity has two local propagating tensor polarizations around a
regular four-dimensional background.  Linearized metric perturbations can be
quantized, and their quanta are gravitons.  At energies well below the Planck
scale, General Relativity is a consistent quantum effective field theory:
higher-curvature operators encode short-distance uncertainty, while
long-distance quantum corrections can be calculated without knowing the
ultraviolet completion \cite{Donoghue}.

Perturbative nonrenormalizability therefore does not mean that quantum gravity
is meaningless.  It means that the Einstein--Hilbert action is not expected
to be the last term in a fundamental ultraviolet description.  This is the
same EFT logic used throughout particle physics, with gravity's very high
suppression scale making the low-energy expansion especially successful.

\subsection{Universal spin-two coupling is a real constraint}

Massless helicity-two consistency strongly constrains how a graviton couples.
Soft-graviton arguments connect Lorentz invariance and factorization to
universal coupling, while nonlinear completion leads toward Einstein
dynamics \cite{WeinbergSoft,Deser}.  A projection-first theory cannot evade
these results by changing vocabulary.  It must recover the same physical pole,
gauge redundancy, conserved source, and self-coupling in the regime where GR
is observed.

\subsection{There is no single established ultraviolet completion}

String theory, loop and spin-foam approaches, asymptotic safety, causal sets,
holography, group-field and tensor models, and other programs attack
different parts of the problem.  Several contain rigorous or highly
developed mathematical sectors; none currently has a unique, experimentally
confirmed description of Planck-scale gravity.  It is more accurate to call
the ultraviolet problem open than to call all these programs failures.

\section{The Projection-First Hypothesis}

\subsection{Geometry as response rather than fundamental substance}

The proposed MTT reading is that effective geometry records the coherent
response to source action and closure strain.  Schematically, after a selected
projection \(\cP\),
\begin{equation}
 S_{\rm eff}[g,\psi]
 =S_{\rm geom}[g]+S_{\rm source}[g,\cP\psi]
 +\Delta S_{\rm proj}[g,\psi],
\label{eq:Seff}
\end{equation}
and the metric response is determined by
\begin{equation}
 \frac{\delta S_{\rm eff}}{\delta g^{\mu\nu}}=0.
\label{eq:response}
\end{equation}
This is compatible with the familiar Einstein equation when
\(S_{\rm geom}\) contains the Einstein--Hilbert term and the projected source
has a conserved stress tensor.  Equations~\eqref{eq:Seff}--\eqref{eq:response}
do not derive either ingredient.  They display the location of the missing
source theorem.

The phrase ``quantized action stresses the medium'' can therefore be made
precise in two stages:
\begin{enumerate}
\item construct quantum source observables and their selected state or
functional; and
\item prove that their variation induces a local, conserved, positive
geometric response with the observed normalization.
\end{enumerate}
Until both stages are supplied, the phrase is an interpretation rather than a
calculation.

\subsection{Emergent and quantum are compatible}

An effective variable may be emergent and still require quantization.  Sound
waves are collective yet quantized as phonons; hydrodynamic variables are
coarse yet fluctuate.  Likewise, if metric perturbations are the low-energy
collective modes of a deeper MTT carrier, their quantum correlations remain
part of the effective theory.  The meaningful distinction is not
``quantum versus geometric'' but:
\[
\text{fundamental source variable}
\quad\hbox{versus}\quad
\text{effective collective variable}.
\]
The projection-first proposal concerns this ontology and the map between
levels.  It does not repeal low-energy graviton physics.

\section{A Concrete Model of Noninvertible Description}

Consider a bipartite Hilbert space
\(\cH_{\rm full}=\cH_{\rm obs}\otimes\cH_{\rm hidden}\).  A restricted
observer assigns
\[
 \rho_{\rm obs}
 =\operatorname{Tr}_{\rm hidden}\rho_{\rm full}.
\]
Many distinct full states have the same reduced density operator, so the
partial trace has no inverse on the full state space.  Entropy can increase
for the restricted description even though the full state evolves
unitarily.  This is an exact and familiar example of effective information
loss without fundamental destruction of information.

The example captures the intuition behind projection-first horizon language:
an observer may lack an inverse map from accessible data to the global state.
It does \emph{not} prove that a gravitational horizon is literally a partial
trace, nor that the hidden factor is an MTT bundle.  A gravity construction
must derive the relevant observable algebra, state restriction, modular
structure, and causal domain rather than import the analogy.

\section{Horizons: A Careful Projection Interpretation}

\subsection{What a horizon is in GR}

An event horizon is a global causal boundary defined relative to the future
of the spacetime.  A freely falling observer need not encounter any local
singularity at a large regular horizon.  The horizon therefore cannot be
identified in general with local failure of admissibility or local breakdown
of physics.

For a restricted exterior observer, however, the inaccessible region changes
which observables can be reconstructed.  Quantum field theory in such
restricted regions, black-hole thermodynamics, and holographic dualities all
make this restriction physically important.  Bekenstein--Hawking entropy
scales with horizon area \cite{Bekenstein,Hawking,GibbonsHawking}, but the area
law alone does not identify the microscopic states or prove a particular MTT
capacity functional.

\subsection{The admissibility hypothesis}

A defensible projection-first hypothesis is:
\begin{quote}
A horizon may mark a boundary of an observer's reconstructible effective
algebra even when the local field equations and freely falling description
remain regular.
\end{quote}
This statement is observer- and algebra-relative.  To promote it to a theorem,
one must construct a restriction map, characterize its kernel, show how the
area law arises, and recover the known Hawking and generalized-entropy
relations.  Saying simply that all interior states project to one exterior
state would be false and far too coarse: exterior observables can retain
substantial information about charges, radiation, correlations, and the
black-hole state.

\section{Singularities: Boundary of a Theory, Not Yet a Projection Theorem}

The singularity theorems establish geodesic incompleteness under stated
causal and energy conditions \cite{HawkingPenrose}.  A singular spacetime need
not contain a point of infinite scalar curvature, and a ``singularity'' is
not ordinarily a point belonging to the manifold.  It is therefore better to
say that classical spacetime evolution becomes inextendible in a specified
sense than to say that every distinction categorically disappears.

Projection-first language can still be useful.  If the effective carrier,
principal symbol, or observable map loses the estimates required for
continuation, then the effective description has reached a mathematical
boundary.  A constructive theory should identify which estimate fails:
\begin{itemize}
\item hyperbolicity or causal stability;
\item boundedness of the projection;
\item positivity of the physical state space;
\item control of curvature and stress-energy;
\item invertibility of a reconstruction map; or
\item convergence of the effective expansion.
\end{itemize}
This replaces a metaphorical ``failure of description'' with a list of
objects that can actually be tested.

\section{Relationship to Other Approaches}

\begin{center}
\begin{tabular}{p{0.22\linewidth}p{0.33\linewidth}p{0.35\linewidth}}
\toprule
Approach & Shared ground & Distinct projection-first question \\
\midrule
Gravity as EFT &
Low-energy gravitons and higher-curvature corrections are valid. &
Can their coefficients and source map descend from one selected upper
carrier? \\
Semiclassical gravity &
Quantum matter sources an effective classical geometry. &
What state-selection and fluctuation terms replace a bare expectation-value
source when that approximation fails? \\
Holography and quantum error correction &
Bulk reconstruction can be observer- or code-subspace dependent
\cite{ADH}. &
Is the MTT projection an explicit encoding/reconstruction map with the same
entropy and locality properties? \\
String theory &
Consistent ten-dimensional quantum systems can yield gravity and gauge
sectors after compactification. &
Can MTT select the compactification, normalized zero modes, and effective
action rather than merely reinterpret them? \\
Loop/spin-network approaches &
Geometric observables can have discrete kinematics. &
Does an MTT finite carrier reproduce the physical constraint algebra and
continuum dynamics, not only a discrete support pattern? \\
\bottomrule
\end{tabular}
\end{center}

The projection-first proposal is most credible when it acts as a precise
bridge among these established structures.  It is least credible when it
declares them unnecessary without reproducing their successful limits.

\section{Constructive Gates for an MTT Quantum-Gravity Sector}

A local admissible QG sector is not complete merely because all loop
integrals in one representation are finite.  At minimum, one must provide:
\begin{enumerate}
\item \textbf{Selected causal carrier.} A Lorentzian principal symbol and
well-posed local evolution, not a positive Gram tensor relabeled as spacetime.
\item \textbf{Physical spin-two sector.} Two helicity-two polarizations, a
positive pole residue, and no unremoved negative-norm modes.
\item \textbf{Gauge consistency.} Diffeomorphism/BRST or an equivalent
constraint complex with anomaly control and a positive physical state space.
\item \textbf{Universal source response.} A normalized coupling to a
conserved stress tensor and the nonlinear self-coupling required in the GR
limit.
\item \textbf{Controlled reduction.} A consistent-truncation equation or an
explicit error bound for discarded modes, including moduli and vector zero
modes.
\item \textbf{Ultraviolet statement.} A regulator-independent construction or
a clearly delimited EFT, with analyticity, unitarity, and causality checked in
the same theory.
\item \textbf{Global gluing.} Compatibility of local sectors on overlaps,
including state, connection, holonomy, and observable maps.
\item \textbf{Empirical dictionary.} A Newton and Planck normalization,
the post-Newtonian and gravitational-wave limits, and at least one
discriminating prediction not used as input.
\end{enumerate}

These gates are interdependent.  For example, a damped Euclidean propagator
does not establish Lorentzian unitarity by itself, and a finite helicity
carrier does not fix the pole residue or stress response.  A successful
construction must make the same selected source pass all gates.

\section{Where MTT Currently Stands}

The current MTT corpus contains three relevant levels:
\begin{itemize}
\item a corrected ten-dimensional action paper that explicitly treats the
metric and Einstein--Hilbert sector as a regime-local ansatz;
\item a controlled coherent-reduction paper that gives the exact discarded
equation and an approximate gap/error theorem; and
\item a successor calculation that closes internal
\(\mathbb Z_{64}\) transverse-traceless support with two real helicity-two
directions and a normalized internal eigenvalue.
\end{itemize}

This is real progress in organizing the gravity bridge.  It establishes that
the finite selected carrier can host the right internal TT support and that a
supplied metric action has a mathematically controlled route to a
four-dimensional GR sector.

The practical content of the result is also limited.  The corpus does not yet
derive the physical Newton constant, Planck scale, complete matter
stress-energy response, positive graviton pole, or a unique ultraviolet
completion from one MTT source.  Earlier constructive-QG packets should
therefore be read as conditional contracts unless their hypotheses have been
closed independently.

\section{What Would Count as Success or Failure}

\subsection{Success}

The program would become substantially stronger if one selected source
operator produced, without importing measured gravity data:
\begin{itemize}
\item the physical Lorentzian spin-two Hessian and its positive residue;
\item a conserved universal matter-response map;
\item \(\kappa_4\) from normalized internal data;
\item a controlled GR infrared limit; and
\item a finite or UV-complete quantum construction with common gauge and
analytic-continuation policies.
\end{itemize}
Agreement among these outputs would turn the projection-first picture from an
interpretation into a constructive theory.

\subsection{Failure}

The proposal would fail in its strong form if no common source can produce
the physical spin-two and stress sectors, if the selected projection violates
unitarity or causality, if extra light modes contradict observation, or if
all numerical agreement requires replaying the measured constants it claims
to derive.  These are useful failure conditions because they keep the
reframing answerable to physics.

\section{Discussion: What This Reframing Means}

The projection-first perspective does not remove quantum gravity.  It
reorganizes the problem.  The quantum object may be a deeper source and
selection structure, while spacetime geometry is its effective collective
response.  The resulting metric must still reproduce every successful
low-energy statement ordinarily described as quantum gravity.

This interpretation gives hope for a specific reason: it permits one common
selection map to organize matter, gauge, and geometric response before their
complicated projected rules are checked separately.  Its risk is equally
clear: unless that map is constructed, the language merely redescribes known
effective physics.  The frontier is therefore an operator and normalization
problem, not a philosophical declaration that gravity should never be
quantized.

\section{Conclusion}

A defensible projection-first reframing says neither that gravity is unreal
nor that metric quantization is meaningless.  It asks whether effective
quantum geometry can be derived as the coherent response of deeper selected
quantum source data.  Restricted descriptions are naturally noninvertible,
so this viewpoint may illuminate why observer-dependent reconstruction,
entropy, and causal boundaries recur in gravitational physics.

The idea remains a research program.  Current MTT results provide internal
helicity-two support and conditional reduction machinery, not full
constructive quantum-gravity closure.  The next decisive advance is a
same-source Lorentzian Hessian and conserved stress-response theorem with
physical normalization.  That object would let the interpretation meet the
successful mathematics of GR, QFT, and existing quantum-gravity approaches
on equal terms.

\begin{thebibliography}{99}

\bibitem{MTTFoundation}
P.~Nero,
\emph{Modal Triplet Theory: Foundations},
revised edition, 2026.

\bibitem{MTTGR}
P.~Nero,
\emph{Controlled Coherent Reduction to Four-Dimensional Einstein Gravity},
third edition, 2026.

\bibitem{MTTAction}
P.~Nero,
\emph{Closure Geometry and a Regime-Local Ten-Dimensional Action Ansatz},
fourth edition, 2026.

\bibitem{Donoghue}
J.~F.~Donoghue,
``General relativity as an effective field theory: The leading quantum
corrections,''
\emph{Physical Review D} \textbf{50} (1994) 3874--3888,
\href{https://doi.org/10.1103/PhysRevD.50.3874}
{doi:10.1103/PhysRevD.50.3874}.

\bibitem{WeinbergSoft}
S.~Weinberg,
``Photons and gravitons in S-matrix theory: Derivation of charge conservation
and equality of gravitational and inertial mass,''
\emph{Physical Review} \textbf{135} (1964) B1049--B1056,
\href{https://doi.org/10.1103/PhysRev.135.B1049}
{doi:10.1103/PhysRev.135.B1049}.

\bibitem{Deser}
S.~Deser,
``Self-interaction and gauge invariance,''
\emph{General Relativity and Gravitation} \textbf{1} (1970) 9--18,
\href{https://doi.org/10.1007/BF00759198}
{doi:10.1007/BF00759198}.

\bibitem{Bekenstein}
J.~D.~Bekenstein,
``Black holes and entropy,''
\emph{Physical Review D} \textbf{7} (1973) 2333--2346,
\href{https://doi.org/10.1103/PhysRevD.7.2333}
{doi:10.1103/PhysRevD.7.2333}.

\bibitem{Hawking}
S.~W.~Hawking,
``Particle creation by black holes,''
\emph{Communications in Mathematical Physics} \textbf{43} (1975) 199--220,
\href{https://doi.org/10.1007/BF02345020}
{doi:10.1007/BF02345020}.

\bibitem{GibbonsHawking}
G.~W.~Gibbons and S.~W.~Hawking,
``Action integrals and partition functions in quantum gravity,''
\emph{Physical Review D} \textbf{15} (1977) 2752--2756,
\href{https://doi.org/10.1103/PhysRevD.15.2752}
{doi:10.1103/PhysRevD.15.2752}.

\bibitem{HawkingPenrose}
S.~W.~Hawking and R.~Penrose,
``The singularities of gravitational collapse and cosmology,''
\emph{Proceedings of the Royal Society A} \textbf{314} (1970) 529--548,
\href{https://doi.org/10.1098/rspa.1970.0021}
{doi:10.1098/rspa.1970.0021}.

\bibitem{ADH}
A.~Almheiri, X.~Dong, and D.~Harlow,
``Bulk locality and quantum error correction in AdS/CFT,''
\emph{Journal of High Energy Physics} \textbf{04} (2015) 163,
\href{https://doi.org/10.1007/JHEP04(2015)163}
{doi:10.1007/JHEP04(2015)163}.

\end{thebibliography}

\end{document}
